\documentclass[11pt]{article}

%------------------------------------------------
% Packages
%------------------------------------------------
\usepackage[a4paper,margin=2.5cm]{geometry}
\usepackage{amsmath,amssymb,amsthm}
\usepackage{mathtools}
\usepackage{enumitem}
\usepackage{hyperref}

\hypersetup{
    colorlinks=true,
    linkcolor=blue,
    citecolor=blue,
    urlcolor=blue
}

%------------------------------------------------
% Title Information
%------------------------------------------------

\title{\bfseries Atomic Gradient Flows: Gradient Flows on Sparse Representations}

\author{
Christian Amend\thanks{Department of Mathematics, University of Twente,
\texttt{christian.amend@utwente.nl}}
\and
Marcello Carioni\thanks{Department of Mathematics, University of Twente,
\texttt{m.c.carioni@utwente.nl}}
\and
Kostantinos Zemas\thanks{Institute of Applied Mathematics, University of Bonn,
\texttt{zemas@iam.uni-bonn.de}}
}

\date{}

%------------------------------------------------

\begin{document}

\maketitle

\vspace{-1em}

\noindent
\textbf{Keywords:}
optimization, inverse problems, gradient flows in metric spaces, sparsity

\vspace{0.5em}


\vspace{1em}

%%%%%%%%%%%%%%%%%%%%%%%%%%%%%%%%%%%%%%%%%%%%%%%%%%
\section*{Abstract}

One of the most popular approaches for optimization problems in the space of measures are the so-called Particle Gradient Flows (PGFs). PGFs restrict the problem to linear combinations of Dirac deltas and then perform a gradient flow in the space of weights and positions \cite{Chizat}. In this talk, I will show how to generalize PGFs to arbitrary spaces and optimization problems:
\begin{align*}
\inf_{u \in X} G(u) \quad \text{where} \quad  G(u) := F(Au) + R(u)
\end{align*}
where $X$ is a Banach space, $F$ is a smooth fidelity term, $A$ is a linear operator and $R$ is a convex regularizer \cite{atomic}.
The crucial idea is how to choose the right notion of particles, which we argue to be the extremal points of the ball of $R$. Such observation is motivated by Krein-Millman theorem that allows to represent any element in the domain of $G$ as the limit of convex combinations of such extremal points. In particular, one considers the restriction of $G$ to \emph{sparse representations}
\begin{align*}
u_{\rm sparse} = \sum_{i=1}^n c_i u_i \quad c_i \in \mathbb{R}_+, \ \ u_i \in {\rm Ext}(\{R \leq 1\}) 
\end{align*} 
and performs a gradient flow of such restriction in the extremal points $u_i$ and coefficients $c_i$, named \emph{Atomic Gradient Flow} (AGF). I will consider analytical properties of such flow, showing its well-posedness as a minimizing movement scheme, its uniqueness and the related contraction estimates. Then, using Choquet’s theorem, I will present how to lift the problem into the Wasserstein space over extremal points and weights, discussing the relations between the lifting of the AGF and the AGF itself. 
Finally, I will show how AGFs can be used to deal with optimization problems with BV-seminorm and Optimal Transport regularizations.  

%%%%%%%%%%%%%%%%%%%%%%%%%%%%%%%%%%%%%%%%%%%%%%%%%%
\vspace{1em}

\noindent
\textbf{Acknowledgements.}
C.A. and M.C.’s research is supported by the NWO-M1 grant Curve Ensemble Gradient Descents for
Sparse Dynamic Problems (OCENW.M.22.302) and NWO-Vidi grant SPARGO (VI.Vidi.243.200).

%%%%%%%%%%%%%%%%%%%%%%%%%%%%%%%%%%%%%%%%%%%%%%%%%%
\vspace{1em}

\begin{thebibliography}{9}

\bibitem{atomic}
C.~Amend, M.~Carioni and K.~Zemas.
\newblock {\em Atomic Gradient Flows: Gradient Flows on Sparse Representations}.
\newblock arXiv preprint arXiv:2603.25675, 2026.

\bibitem{Chizat}
L. Chizát, F. Bach.
\newblock {\em On the Global Convergence of Gradient Descent for Over-parameterized Models using Optimal Transport}.
\newblock Neurips, 2018.

\end{thebibliography}

\end{document}